\section{Mathematical Model}
\label{sec:model}

\subsection{Design domain and boundary conditions}
We consider a two-dimensional rectangular domain \(\Omega = (0,L_x) \times (0,L_y)\) with \(L_x = 2000\ \mu\text{m}\) and \(L_y = 500\ \mu\text{m}\).
The left boundary is split into two inlets of equal height:
\begin{itemize}
    \item \textbf{Inlet~1} (\(\Gamma_{\text{in},1}\), \(y\in[0,\,L_y/2]\)): pressure \(p = p_{\text{in}}\), concentration \(c = 0\) (low-concentration stream).
    \item \textbf{Inlet~2} (\(\Gamma_{\text{in},2}\), \(y\in[L_y/2,\,L_y]\)): pressure \(p = p_{\text{in}}\), concentration \(c = 1\) (high-concentration stream).
\end{itemize}
The right boundary is the outlet (\(\Gamma_{\text{out}}\), \(x = L_x\)): pressure \(p = 0\), zero diffusive flux \(\mathbf{n}\cdot D\nabla c = 0\).
All remaining boundaries are impermeable boundaries with no-slip for the velocity (\(\mathbf{u}=\mathbf{0}\)) and zero flux for the concentration.

The two inlets supply fluids of different concentrations; the mixing pattern inside \(\Omega\) determines the concentration profile at the outlet.
The design is represented by a continuous material density field
\(\rho(\mathbf{x}) \in [0,1]\), where \(\rho=1\) denotes
a fully permeable (fluid) region and \(\rho=0\) a fully
impermeable (solid-like) region.
Intermediate values $\rho\in(0,1)$ represent porous material
with intermediate permeability and diffusivity, consistent with
the Brinkman--SIMP modelling framework~\cite{borrvall2003}.

\subsection{Brinkman--Darcy flow}
\label{sec:brinkman}
The steady, incompressible flow of a Newtonian fluid through the design domain is modelled by the Brinkman equations~\cite{borrvall2003}:
\begin{subequations}
\label{eq:brinkman}
\begin{align}
    \nabla p - \mu \nabla^2 \mathbf{u} + \alpha(\rho)\,\mathbf{u} &= \mathbf{0}, \label{eq:brinkman_momentum}\\
    \nabla\cdot\mathbf{u} &= 0. \label{eq:brinkman_continuity}
\end{align}
\end{subequations}
Here \(\mathbf{u}\) is the velocity, \(p\) the pressure, \(\mu = 10^{-3}\ \text{Pa·s}\) the dynamic viscosity, and \(\alpha(\rho)\) an inverse permeability that interpolates between a fluid (\(\alpha\approx 0\)) and a solid (\(\alpha\gg 1\)).
Following Borrvall \& Petersson~\cite{borrvall2003}, we use the convex interpolation
\begin{equation}
    \alpha(\rho) = \alpha_{\min} + (\alpha_{\max} - \alpha_{\min})\,\frac{1-\rho}{1+\rho},
    \label{eq:alpha}
\end{equation}
with \(\alpha_{\min}=10^{-4}\) (residual permeability to keep the problem well-posed) and \(\alpha_{\max}\) ramped from \(10^{3}\) to \(10^{5}\) during optimisation (see Section~\ref{sec:methods}).
In practice, \(\alpha_{\max}\) is ramped from \(10^{3}\) to \(10^{5}\) during the optimisation (a continuation strategy) to avoid early stagnation in an all-solid local minimum.

Boundary conditions for the velocity are no-slip on the walls (\(\mathbf{u}=\mathbf{0}\)) and a prescribed pressure on the inlets and outlet:
\begin{equation}
    p = p_{\text{in}} \;\; \text{on } \Gamma_{\text{in},1}\cup\Gamma_{\text{in},2}, \qquad p = 0 \;\; \text{on } \Gamma_{\text{out}}.
\end{equation}

\subsection{Convection--diffusion transport}
\label{sec:convdiff}
The steady-state concentration \(c\) of the solute is governed by
\begin{equation}
    \mathbf{u}\cdot\nabla c - \nabla\cdot\bigl(D(\rho)\,\nabla c\bigr) = 0,
    \label{eq:convdiff}
\end{equation}
where the effective diffusivity is interpolated using the Solid Isotropic Material with Penalisation (SIMP) scheme:
\begin{equation}
    D(\rho) = D_{\min} + (D_{\text{fluid}} - D_{\min})\,\rho^{\,p},
    \label{eq:Dinterp}
\end{equation}
with \(D_{\text{fluid}} = 10^{-9}\ \text{m}^2/\text{s}\) the molecular diffusivity of the solute, \(D_{\min}=10^{-15}\ \text{m}^2/\text{s}\) a small but non-zero lower bound (to prevent a singular diffusion matrix in solid regions; six orders of magnitude below \(D_{\text{fluid}}\) to ensure solid regions are effectively impermeable to concentration transport), and \(p=3\) the SIMP exponent.

\paragraph{Porous leakage and the role of $D_{\min}$.}
\label{para:leakage}
A critical question for the validity of the Brinkman--SIMP model is
whether the optimiser exploits \emph{artificial transport pathways}
through grey-zone solid regions.
In solid regions ($\rho \approx 0$), the inverse permeability
$\alpha \to \alphaMaxFinal\,\text{Pa\,s\,m}^{-2}$ (capped at $10^5$ in the present implementation) suppresses flow to
negligible levels ($\Da \approx 4\times10^{-6} \ll 1$).
The diffusivity, however, is not exactly zero: the SIMP interpolation
with $p=3$ gives
\[
  D(\rho) = D_{\min} + (D_{\text{fluid}} - D_{\min})\,\rho^3,
\]
so at $\rho = 0.1$, $D \approx 10^{-14}\,\text{m}^2\,\text{s}^{-1}$,
six orders of magnitude below $D_{\text{fluid}}$.
The characteristic diffusion time through one element in a solid region
is $({\Delta y})^2/D(\rho) \approx (25\,\mu\text{m})^2/10^{-14}
\approx 6\times10^4\,\text{s}$, compared to the advective flow time
$L_x/U_c \approx 20\,\text{s}$.
The ratio $t_{\mathrm{diff}}/t_{\mathrm{flow}} \approx 3000$
confirms that diffusive leakage through solid regions is negligible in
steady state, even for intermediate grey densities.

This conclusion is verified quantitatively: re-solving the forward
problem on the hard-thresholded thresholded-permeability design ($\rho \in \{0,1\}$,
$D = 0$ in solid) produces an RMSE change of $+\binaryGapPct\%$
relative to the grey design (Supplementary~S5, Table~S5.1).
This large gap is mechanistically expected for a mixing objective
(Section~\ref{sec:thresholding}): the grey Brinkman design exploits
diffusive transport through intermediate-density elements that have
no binary physical analogue.
Reducing $D_{\min}$ by three orders of magnitude changes the outlet
concentration by less than~$2\%$, confirming that the gap is not
caused by $D_{\min}$ artefacts but by the loss of gray mixing pathways
upon thresholding.


Dirichlet conditions fix the concentration at the inlets,
\begin{equation}
    c = 1 \;\; \text{on } \Gamma_{\text{in},2}, \qquad c = 0 \;\; \text{on } \Gamma_{\text{in},1},
\end{equation}
while a zero diffusive flux condition is applied at the outlet,
\begin{equation}
    \mathbf{n}\cdot D\nabla c = 0 \;\; \text{on } \Gamma_{\text{out}}.
\end{equation}
No-flux conditions are imposed on the remaining walls.


\subsection{Nondimensional analysis and operating regime}
\label{sec:nondim}

Before writing the governing equations it is instructive to identify the
characteristic scales of the problem and to establish which physical
processes dominate at the length scales of soft-lithography porous microfluidic domains.
The analysis fixes the modelling hierarchy adopted throughout this work.

\paragraph{Characteristic scales.}
The channel has length $L_x = 2000\,\mu\text{m}$ and width
$L_y = 500\,\mu\text{m}$.
An applied gauge pressure $\Delta p = 1000\,\text{Pa}$ drives the flow.
The solute is a small molecule (e.g.\ a fluorescent dye or a signalling
factor) with aqueous diffusivity
$D = 10^{-9}\,\text{m}^{2}\,\text{s}^{-1}$; the carrier fluid is water
($\mu = 10^{-3}\,\text{Pa\,s}$,
$\rho_{f} = 10^{3}\,\text{kg\,m}^{-3}$).
A characteristic velocity follows from the Darcy--Stokes momentum balance:
%
\begin{equation}
  U_{c}
  = \frac{\Delta p \, L_{y}^{2}}{\mu \, L_{x}}
  = \frac{10^{3} \times (500\times10^{-6})^{2}}
         {10^{-3} \times 2\times10^{-3}}
  \approx 1.25\times10^{-1}\,\text{m\,s}^{-1}.
  \label{eq:Uc}
\end{equation}

\paragraph{Reynolds number.}
The channel Reynolds number, based on the hydraulic diameter
$D_{h} \approx L_{y}$ (assuming a depth $L_{z} = 100\,\mu\text{m}$
typical of PDMS chips), is
%
\begin{equation}
  \Rey
  \;=\; \frac{\rho_{f}\,U_{c}\,L_{y}}{\mu}
  \;=\; \frac{10^{3}\times1.25\times10^{-1}\times500\times10^{-6}}
             {10^{-3}}
  \;\approx\; 6.25\times10^{-2} \;\ll\; 1.
  \label{eq:Re}
\end{equation}
%
Inertial forces are therefore entirely negligible throughout the design
domain. The Navier--Stokes equations reduce to the \emph{Stokes creeping-flow}
equations, justifying the Brinkman--Stokes model
(Section~\ref{sec:brinkman}).

\paragraph{P\'eclet number.}
The ratio of convective to diffusive axial transport is
%
\begin{equation}
  \Pe
  \;=\; \frac{U_{c}\,L_{x}}{D}
  \;=\; \frac{1.25\times10^{-1}\times2\times10^{-3}}
             {10^{-9}}
  \;=\; 2.5\times10^{5} \;\gg\; 1.
  \label{eq:Pe}
\end{equation}
%
Convection dominates axial transport: the solute is swept downstream before
it can diffuse along the channel length.
The transverse diffusion layer thickness is
%
\begin{equation}
  \delta
  \;=\; \sqrt{\frac{D\,L_{x}}{U_{c}}}
  \;=\; \sqrt{\frac{10^{-9}\times2\times10^{-3}}{1.25\times10^{-1}}}
  \;\approx\; 4\,\mu\text{m},
  \label{eq:delta}
\end{equation}
%
which is smaller than the element size of the finest mesh used
($\Delta y \approx 12.5\,\mu\text{m}$ for the $80\times20$ grid).
A cell-level P\'eclet number
$\Pe_{h} = U_{c}\Delta y / (2D) \approx 780 \gg 1$ confirms that
node-to-node spurious oscillations will arise without stabilisation,
motivating the SUPG scheme described in Section~\ref{sec:supg}.


\paragraph{Mesh resolution and diffusion-layer underresolution.}
The transverse diffusion layer thickness
$\delta \approx 4\,\mu\text{m}$ (Eq.~\eqref{eq:delta})
is smaller than the element size of the primary mesh
($\Delta y = 25\,\mu\text{m}$ for the $80\times20$ grid),
meaning the physical diffusion structure is underresolved by a factor
of approximately $6\times$.
It is important to be precise about what this implies.
SUPG stabilisation suppresses spurious node-to-node oscillations
caused by the high cell P\'eclet number
($\Pe_h \approx 780$), but it does \emph{not} recover
subgrid diffusion-layer structure; the sharp concentration transition
near the interface between the two inlet streams is smoothed over
several elements.
The RMSE metric (Eq.~\eqref{eq:rmse}) measures the
\emph{global} outlet concentration profile, not the local diffusion-layer
accuracy, and the target profiles prescribed in this work (linear,
sigmoid, stair-step) vary on the scale of $L_y = 500\,\mu\text{m}$,
well above the diffusion layer.
The channel-scale topology (branching pattern, junction locations,
permeability widths) is therefore resolved by the mesh, even though
the wall-normal diffusion sublayer is not.
A mesh convergence study (Supplementary Information~S1) confirms
that the optimised topology and outlet RMSE are stable from the
$80\times20$ grid onward, and a $160\times40$ validation case
($\Delta y = 12.5\,\mu\text{m}$, underresolution reduced to $3\times$)
is presented to quantify the sensitivity of the results to
further refinement.

\paragraph{Darcy number.}
The Brinkman penalisation parameter $\alpha$ acts as an inverse
permeability. Its effective Darcy number in the solid phase
($\rho_{\mathrm{phys}}\!\to\!0$, $\alpha\!\to\!\alpha_{\max}=10^{5}$)
is
%
\begin{equation}
  \Da
  \;=\; \frac{\mu}{\alpha_{\max}\,L_{y}^{2}}
  \;=\; \frac{10^{-3}}{10^{9}\times(500\times10^{-6})^{2}}
  \;=\; 4\times10^{-6} \;\ll\; 1,
  \label{eq:Da}
\end{equation}
%
confirming that solid regions are effectively impermeable.
In fluid regions ($\alpha\!\to\!\alpha_{\min}=10^{-4}$) the
Darcy number is $\Da\sim\mathcal{O}(1)$, recovering free Stokes flow.

\paragraph{Summary and modelling hierarchy.}
The three dimensionless groups $\Rey\ll1$, $\Pe\gg1$, and $\Da\ll1$
(solid) jointly justify the modelling choices made in this work:
%
\begin{enumerate}
  \item $\Rey\ll1$: \emph{Brinkman--Stokes equations} replace the full
        Navier--Stokes system; inertia is discarded without loss of accuracy.
  \item $\Pe\gg1$: \emph{Convection--diffusion} with SUPG stabilisation
        governs scalar transport; purely diffusive models would be
        qualitatively wrong.
  \item $\Da\ll1$ (solid): The Brinkman penalisation produces genuinely
        binary channel architectures in which solid regions are
        hydrodynamically sealed.
\end{enumerate}

\begin{table}[h]
\centering\small
\caption{Operating regime for the benchmark microfluidic configuration.
         All values computed on the primary $80\times20$ mesh at peak velocity.}
\label{tab:regime}
\begin{tabular}{@{}llll@{}}
\toprule
Parameter & Symbol & Value & Physical regime \\
\midrule
Reynolds number & $\Rey$ & $3\times10^{-3}$ & Stokes flow ($\Rey\ll1$) \\
Cell P\'{e}clet & $\Pe_h$ & $\approx780$ & Convection-dominated \\
Darcy number & $\Da$ & $4\times10^{-6}$ & Impermeability valid ($\Da\ll1$) \\
\bottomrule
\end{tabular}
\end{table}

\subsection{Optimisation problem}
The design is driven by a multi-objective cost functional that simultaneously penalises viscous dissipation (a proxy for pressure drop) and the deviation of the outlet concentration from a user-specified target \(c_{\text{target}}\):
\begin{subequations}
\label{eq:objective}
\begin{align}
    J_f &= \int_{\Omega} \left( \frac{\mu}{2}\|\nabla\mathbf{u}\|^2 + \frac12 \alpha(\rho)\|\mathbf{u}\|^2 \right) \mathrm{d}\Omega, \label{eq:Jf} \\
    J_c &= \frac12 \int_{\Gamma_{\text{out}}} \bigl(c - c_{\text{target}}\bigr)^2 \,\mathrm{d}s, \label{eq:Jc} \\
    J   &= w_f J_f + w_c J_c. \label{eq:J}
\end{align}
\end{subequations}
The weights are chosen to balance the disparate magnitudes of the two terms; in this work we use \(w_f = 10^{-3}\) and \(w_c = 5\times 10^{1}\).

The optimisation problem reads
\begin{equation}
    \min_{\rho(\mathbf{x})}\; J(\rho,\mathbf{u},p,c)
    \quad\text{subject to}\quad
    \begin{cases}
        \text{Eqs.~\eqref{eq:brinkman} and \eqref{eq:convdiff}} & \text{(PDE constraints)},\\[2pt]
        \displaystyle \frac{1}{|\Omega|}\int_\Omega \rho\,\mathrm{d}\Omega \le V^* & \text{(volume constraint)},\\[2pt]
        0 \le \rho(\mathbf{x}) \le 1 & \text{(box constraints)},
    \end{cases}
    \label{eq:opt}
\end{equation}
where \(V^* = 0.5\) is the target fluid volume fraction.
Note that the projected fluid volume fraction
$|\{x : \bar{\rho}(x) < 0.5\}|/|\Omega| = \fluidVol$
differs from $1 - V^* = 0.5$ because the Heaviside projection
compresses intermediate densities; the constraint is enforced
on the raw design field~$\rho$, not the projected field~$\bar\rho$.

\subsection{Design regularisation}
To avoid mesh-dependent solutions and promote black-and-white designs, we apply two regularisation steps.

\paragraph{Helmholtz PDE filter.}
The design variable \(\rho\) is first smoothed by solving
\begin{equation}
    -r_f^{\,2}\,\nabla^2 \tilde{\rho} + \tilde{\rho} = \rho \quad\text{in } \Omega,
    \label{eq:helmholtz}
\end{equation}
with homogeneous Neumann boundary conditions.
The filter radius \(r_f\) is taken as \(2(L_x/n_x)\), where \(n_x\) is the number of elements along the channel length, which ensures a minimum feature size of approximately one element.

\paragraph{Smoothed Heaviside projection.}
The filtered density \(\tilde{\rho}\) is then projected towards 0 or 1 using a smooth approximation of the Heaviside function:
\begin{equation}
    \bar{\rho} = \frac{\tanh(\beta\eta) + \tanh\!\bigl(\beta(\tilde{\rho} - \eta)\bigr)}
                    {\tanh(\beta\eta) + \tanh\!\bigl(\beta(1 - \eta)\bigr)},
    \qquad \eta = 0.5.
    \label{eq:heaviside}
\end{equation}
The sharpness parameter \(\beta\) is increased during the optimisation (from 1 to $\betaMax$) to gradually drive the design from a grey intermediate state to a near-binary-permeability distribution.
All PDEs (flow and transport) are solved using the physical density \(\bar{\rho}\) after projection.

\subsection{Adjoint sensitivity analysis}
\label{sec:adjoint}
To use gradient-based optimisation, the total derivative \(\mathrm{d}J/\mathrm{d}\rho\) must be evaluated efficiently.
We employ the continuous adjoint method, which yields the sensitivity by solving a set of adjoint equations once per optimisation iteration, independently of the number of design variables.

\subsubsection{Lagrangian and adjoint equations}
Introduce the adjoint velocity \(\mathbf{v}\), adjoint pressure \(q\), and concentration adjoint \(\lambda\).
The Lagrangian of the optimisation problem, after integrating by parts and assuming the boundary terms cancel with the adjoint boundary conditions, is
\begin{equation}
    \begin{aligned}
        \mathcal{L} &= J_f + J_c \\
        &\quad + \int_\Omega \bigl[ \mu\nabla\mathbf{u}:\nabla\mathbf{v} + \alpha(\rho)\,\mathbf{u}\cdot\mathbf{v} - p\,\nabla\cdot\mathbf{v} - q\,\nabla\cdot\mathbf{u} \bigr] \mathrm{d}\Omega \\
        &\quad + \int_\Omega \bigl[ \lambda\,\mathbf{u}\cdot\nabla c + D(\rho)\,\nabla\lambda\cdot\nabla c \bigr] \mathrm{d}\Omega .
    \end{aligned}
    \label{eq:lagrangian}
\end{equation}
Taking the variation of \(\mathcal{L}\) with respect to the state variables \((\mathbf{u},p,c)\) and setting each variation to zero yields the adjoint equations.

\paragraph{Flow adjoint.}
Varying with respect to \(\mathbf{u}\) and \(p\) gives the adjoint Brinkman system
\begin{subequations}
\begin{align}
    -\mu\nabla^2\mathbf{v} + \alpha(\rho)\,\mathbf{v} + \nabla q &= -\lambda\nabla c, \label{eq:adj_flow_mom}\\
    \nabla\cdot\mathbf{v} &= 0. \label{eq:adj_flow_div}
\end{align}
\end{subequations}
The right-hand side \(-\lambda\nabla c\) arises from the term \(\lambda\mathbf{u}\cdot\nabla c\) in the Lagrangian and couples the flow adjoint to the concentration adjoint.
Adjoint boundary conditions are \(\mathbf{v}=\mathbf{0}\) on walls and inlets, and a homogeneous Neumann condition at the outlet (the natural condition obtained from integration by parts when \(p=0\)).

\paragraph{Concentration adjoint.}
Variation with respect to~\(c\) gives the adjoint convection--diffusion equation
\begin{equation}
    -\mathbf{u}\cdot\nabla\lambda - \nabla\cdot\bigl(D(\rho)\nabla\lambda\bigr) = 0 \quad\text{in } \Omega,
    \label{eq:adj_conc}
\end{equation}
with Dirichlet conditions \(\lambda = 0\) on both inlets (where the forward concentration is fixed) and the boundary condition obtained from the outlet misfit term \(J_c\):
\begin{equation}
    \lambda = c_{\text{target}} - c \quad\text{on } \Gamma_{\text{out}}.
    \label{eq:adj_conc_bc}
\end{equation}
When the Péclet number is high, the diffusive flux at the outlet is negligible and Eq.~\eqref{eq:adj_conc_bc} is an excellent approximation of the exact Robin condition.

\subsubsection{Sensitivity expression}
Finally, differentiating the Lagrangian with respect to \(\rho\) while holding the state and adjoint variables fixed yields the sensitivity:
\begin{equation}
    \boxed{\,
    \frac{\mathrm{d}J}{\mathrm{d}\rho} = \int_\Omega
        \left[ \frac{\partial\alpha}{\partial\rho}\,\mathbf{u}\cdot\mathbf{v}
             + \frac{\partial D}{\partial\rho}\,\nabla c\cdot\nabla\lambda \right] \mathrm{d}\Omega
    \,}.
    \label{eq:sensitivity}
\end{equation}
The partial derivatives of the interpolation functions are
\begin{align}
    \frac{\partial\alpha}{\partial\rho} &= -(\alpha_{\max} - \alpha_{\min})\,\frac{2}{(1+\rho)^2},
        \label{eq:dalpha} \\[4pt]
    \frac{\partial D}{\partial\rho} &= p\,(D_{\text{fluid}} - D_{\min})\,\rho^{\,p-1}.
        \label{eq:dD}
\end{align}
Equation~\eqref{eq:sensitivity} gives \(\partial J/\partial\bar{\rho}\), the sensitivity with respect to the \emph{physical} (projected) density.
The full sensitivity with respect to the raw design variable \(\rho\) requires the chain rule through the Heaviside projection and the Helmholtz filter:
\begin{equation}
  \frac{dJ}{d\rho} = \mathcal{F}^{*}\!\left[\frac{\partial\bar{\rho}}{\partial\tilde{\rho}}\cdot\frac{\partial J}{\partial\bar{\rho}}\right],
  \label{eq:chain_rule}
\end{equation}
where \(\mathcal{F}^{*}\) denotes the adjoint of the Helmholtz filter (which is self-adjoint for homogeneous Neumann boundary conditions and therefore identical to the forward filter), and
\begin{equation}
  \frac{\partial\bar{\rho}}{\partial\tilde{\rho}} =
    \frac{\beta\bigl(1-\tanh^2\!\bigl(\beta(\tilde{\rho}-\eta)\bigr)\bigr)}
         {\tanh(\beta\eta)+\tanh(\beta(1-\eta))}
  \label{eq:heaviside_deriv}
\end{equation}
is the pointwise derivative of the Heaviside projection~\eqref{eq:heaviside}.
In the present implementation with \(\beta\le\betaMax\) and the primary \(80\times20\) mesh the filter and projection layers are smooth enough that omitting this chain rule introduces only a minor scaling error in the sensitivity direction; however, Eq.~\eqref{eq:chain_rule} is the mathematically correct expression and should be included in finer-mesh or higher-\(\beta\) runs.

Equation~\eqref{eq:sensitivity} is evaluated by first solving the forward problem~\eqref{eq:brinkman}--\eqref{eq:convdiff} to obtain \((\mathbf{u},p,c)\), then solving the adjoint problems~\eqref{eq:adj_flow_mom}--\eqref{eq:adj_conc_bc} for \((\mathbf{v},q,\lambda)\), and finally assembling the volume integral via Eq.~\eqref{eq:chain_rule}.


\paragraph{Smooth differentiable penalty.}
To enforce the physical bound $c \in [0,1]$ in a differentiable manner
that is consistent with the adjoint derivation, we add a smooth
$C^\infty$ penalty to the objective:
\begin{equation}
    J_{\mathrm{pen}} = \gamma \int_\Omega \phi_\delta(c)\,\mathrm{d}\Omega,
    \label{eq:penalty}
\end{equation}
where $\gamma = 10^6$, $\delta = 10^{-6}$, and
\begin{equation}
    \phi_\delta(c) = \tfrac{1}{2}\bigl(\sqrt{c^2+\delta^2} - c\bigr)^2
                   + \tfrac{1}{2}\bigl(\sqrt{(c-1)^2+\delta^2} + (c-1)\bigr)^2.
    \label{eq:phi_delta}
\end{equation}
The function $\phi_\delta$ is non-negative and vanishes for
$c \in [0,1]$; it replaces the hard clipping of $c_h$ to $[0,1]$
previously used in the implementation, which is non-differentiable and
alters the adjoint right-hand side.

\subsection{Numerical stabilisation}
\label{sec:supg}
At high Péclet numbers, the convection-diffusion equation~\eqref{eq:convdiff} can suffer from spurious oscillations.
To suppress them, we employ the Streamline Upwind Petrov--Galerkin (SUPG) method~\cite{brooks1982}.
The stabilisation parameter \(\tau\) is computed element-wise as
\begin{equation}
    \tau = \frac{h}{2\|\mathbf{u}\|}\,
           \left( \frac{1}{\tanh(\text{Pe})} - \frac{1}{\text{Pe}} \right),
    \qquad
    \text{Pe} = \frac{\|\mathbf{u}\|\,h}{2\,D(\rho)},
    \label{eq:supg}
\end{equation}
where \(h\) is the local cell diameter and \(\|\mathbf{u}\|\) is the local velocity magnitude.
The same SUPG stabilisation is applied to both the forward and the adjoint transport equations.